\documentclass[anon]{colt2026}

\usepackage[utf8]{inputenc}
\usepackage[T1]{fontenc}
\usepackage{times}
\usepackage{bm}
\usepackage{mathtools}
\usepackage{cleveref}

\newcommand{\RR}{\mathbb{R}}
\newcommand{\CC}{\mathbb{C}}
\newcommand{\EE}{\mathbb{E}}
\newcommand{\PP}{\mathbb{P}}
\newcommand{\Id}{\mathrm{Id}}
\newcommand{\Diag}{\mathrm{Diag}}
\newcommand{\tr}{\mathrm{Tr}}
\newcommand{\od}{\odot}
\newcommand{\ip}[2]{\left\langle #1,#2\right\rangle}
\newcommand{\Norm}[1]{\left\|#1\right\|}
\newcommand{\Cplus}{\{z\in\CC:\Im z>0\}}

\newtheorem{thm}{Theorem}
\newtheorem{lem}{Lemma}
\newtheorem{prop}{Proposition}
\newtheorem{cor}{Corollary}
\newtheorem{defn}{Definition}
\newtheorem{assumption}{Assumption}
\newtheorem{rmk}{Remark}

\title[Deep NTK in quadratic scaling]{Eigenvalue distribution of deep Neural Tangent Kernels in the quadratic scaling}

\begin{document}
\maketitle

\begin{abstract}
We extend the quadratic-scaling spectral analysis of the Neural Tangent Kernel (NTK) from two-layer networks to fixed-depth multilayer perceptrons with linear width ratios.
In the regime $n\asymp d_0 d_1$ with $d_\ell/d_{\ell-1}\to \gamma_{2,\ell}$, each layerwise NTK contribution is a Hadamard product of two low-rank Gram matrices but has nontrivial bulk spectrum.
Our approach reuses the proof spine of \citet{benigni2025eigenvaluedistributionneuraltangent}: screening low-rank terms, rewriting Hadamard products as Gram matrices in a tensor-product feature space, and applying a Bai--Zhou type Marchenko--Pastur map for independent-row sample covariances.
The only genuinely new ingredient is a deep Gaussian-equivalence step: finite-width cumulant expansions \citep{guillen2025finitewidthneuraltangentkernels} imply that the relevant fourth-order tensors match their Gaussian counterparts up to $O(1/d)$ deterministic deformations, which are precisely the perturbations carried by the \citet{benigni2025eigenvaluedistributionneuraltangent} machinery.
\end{abstract}

\begin{keywords}
Neural Tangent Kernel, random matrix theory, quadratic scaling, deep networks, Marchenko--Pastur, finite-width corrections
\end{keywords}

\section{Introduction}
Understanding which functions are learned quickly by gradient descent in neural networks is controlled by the spectrum of the NTK at initialization.
In the two-layer case, \citet{benigni2025eigenvaluedistributionneuraltangent} identified the limiting eigenvalue distribution of the NTK in a \emph{quadratic} high-dimensional regime, in which the sample size is of the same order as the product of the two relevant widths.
The purpose of this paper is to generalize this analysis to fixed-depth multilayer perceptrons with \emph{linear} width ratios, where every layer output remains linear in the input dimension.

Our contribution is conceptual and technical.
Conceptually, we show that the deep NTK admits the same tensor-product Gram representation that underlies the two-layer analysis.
Technically, we isolate the only new obstruction in depth: the ``diagonal weight'' that appears in the two-layer Hadamard product becomes a deep backpropagated signal, and hence is not a priori independent of the same-layer features.
We resolve this by importing finite-width cumulant expansions (Feynman-diagram methods) to justify a Gaussian-equivalence/replacement step that restores the BP setting up to $O(1/d)$ deterministic tensor deformations.

\paragraph{Roadmap.}
In \Cref{sec:model} we define the deep network, the scaling regime, and the NT’K decomposition into layerwise Hadamard products.
In \Cref{sec:prelim} we record the Marchenko--Pastur map viewpoint and a rank-based Stieltjes transform estimate used to discard low-rank terms.
In \Cref{sec:decomp} we tensorize the deep NTK into a single Gram matrix with independent rows.
In \Cref{sec:main} we state the main theorem and outline the proof by matching the dependency graph of \citet{benigni2025eigenvaluedistributionneuraltangent}.

\section{Preliminaries: Stieltjes transforms and the MP map}
\label{sec:prelim}

\begin{defn}[Empirical spectral distribution and Stieltjes transform]
For a real symmetric matrix $A\in\RR^{n\times n}$ with eigenvalues $\lambda_1,\dots,\lambda_n$, the empirical spectral distribution is
\[
  \mu_A \coloneqq \frac{1}{n}\sum_{i=1}^n \delta_{\lambda_i}.
\]
Its Stieltjes transform is $m_A(z)\coloneqq \int \frac{1}{x-z}\,\mu_A(\mathrm{d}x)=\frac{1}{n}\tr(A-z\Id)^{-1}$ for $z\in\Cplus$.
\end{defn}

\begin{lem}[Rank bound for Stieltjes transforms]\label{lem:rank-st}
Let $A,B$ be real symmetric $n\times n$ matrices.
Then for any $z\in\Cplus$,
\[
  \left|
  \frac{1}{n}\tr(A-z\Id)^{-1}-\frac{1}{n}\tr(B-z\Id)^{-1}
  \right|
  \leq \frac{\pi}{n\,\Im z}\,\mathrm{rank}(A-B).
\]
\end{lem}

\begin{defn}[Marchenko--Pastur map]\label{def:mp-map}
Let $\nu$ be a probability measure on $[0,\infty)$ and $\gamma>0$.
The \emph{Marchenko--Pastur map} $\mu_{\mathrm{MP}}^{\gamma}\boxtimes \nu$ is the free multiplicative convolution of $\nu$ with the Marchenko--Pastur law of shape $\gamma$.
\end{defn}

\section{Model and scaling}
\label{sec:model}
Fix a depth $L\geq 1$.
Let $X^{(0)}\in\RR^{n\times d_0}$ be the data matrix with i.i.d.\ entries (centered, unit variance, finite fourth moment).
For $\ell=1,\dots,L$, let $W^{(\ell)}\in\RR^{d_{\ell-1}\times d_\ell}$ have i.i.d.\ $\mathcal{N}(0,1)$ entries, independent across layers and independent of $X^{(0)}$.
Let $a\in\RR^{d_L}$ have i.i.d.\ entries, centered and bounded, independent of all $W^{(\ell)}$.
Define
\[
  Z^{(\ell)} \coloneqq \frac{1}{\sqrt{d_{\ell-1}}}X^{(\ell-1)}W^{(\ell)},\qquad
  X^{(\ell)} \coloneqq \sigma(Z^{(\ell)}),
\]
with $\sigma$ applied entrywise, and consider the scalar-output network
\[
  f(X^{(0)}) \coloneqq \frac{1}{\sqrt{d_L}}X^{(L)}a\in\RR^n.
\]
Let $\phi\coloneqq \sigma'$ (defined Lebesgue-a.e.).
Define the backpropagated signals $\Delta^{(\ell)}\in\RR^{n\times d_\ell}$ by
\[
  \Delta^{(L)} \coloneqq \phi(Z^{(L)})\odot (\mathbf{1}\, a^\top),
  \qquad
  \Delta^{(\ell)} \coloneqq \phi(Z^{(\ell)})\odot \frac{1}{\sqrt{d_\ell}}\Delta^{(\ell+1)}W^{(\ell+1)\top}.
\]

\begin{assumption}[Quadratic scaling]\label{ass:scaling}
There exist constants $\gamma_1>0$ and $\gamma_{2,\ell}>0$ such that as $n\to\infty$,
\[
  \frac{n}{d_0d_1}\to \gamma_1,
  \qquad
  \frac{d_\ell}{d_{\ell-1}}\to \gamma_{2,\ell},
  \qquad \ell=1,\dots,L.
\]
Set $m\coloneqq \sum_{\ell=1}^L d_{\ell-1}d_\ell$ and assume $n/m\to \gamma\in(0,\infty)$.
\end{assumption}

\begin{assumption}[Activation assumptions]\label{ass:activation}
The activation $\sigma$ is locally Lipschitz and differentiable Lebesgue-a.e.
Its derivative $\phi=\sigma'$ is pseudo-Lipschitz and has finite moments under $\mathcal{N}(0,1)$.
This includes smooth activations and can be extended to ReLU by approximation since $\PP(\mathcal{N}(0,1)=0)=0$.
\end{assumption}

\section{Deep NTK decomposition and tensorization}
\label{sec:decomp}
Let $\Theta^{(L)}$ be the empirical NTK at initialization,
\[
  \Theta^{(L)}\coloneqq \nabla_\theta f(X^{(0)})\,\nabla_\theta f(X^{(0)})^\top.
\]
It splits as a low-rank conjugate-kernel term and $L$ Hadamard terms,
\[
  \Theta^{(L)} = K^{\mathrm{CK}} + \sum_{\ell=1}^L K^{(\ell)},
  \qquad
  K^{\mathrm{CK}}=\frac{1}{d_L}X^{(L)}X^{(L)\top},
\]
where for each $\ell$,
\[
  K^{(\ell)} \coloneqq
  \frac{1}{d_{\ell-1}}X^{(\ell-1)}X^{(\ell-1)\top}
  \odot
  \frac{1}{d_\ell}\Delta^{(\ell)}\Delta^{(\ell)\top}.
\]
Since $\mathrm{rank}(K^{\mathrm{CK}})\le d_L$ and $d_L/n\to 0$ in \Cref{ass:scaling}, the bulk ESD is governed by $\sum_{\ell}K^{(\ell)}$.

Define the tensor-product features
\[
  \omega_i^{(\ell)} \coloneqq x_i^{(\ell-1)}\otimes \delta^{(\ell)}(x_i)\in \RR^{d_{\ell-1}d_\ell},
\]
where $x_i^{(\ell-1)}$ is the $i$-th row of $X^{(\ell-1)}$ and $\delta^{(\ell)}(x_i)$ is the $i$-th row of $\Delta^{(\ell)}$.
Then
\[
  K^{(\ell)}_{ij}=\frac{1}{d_{\ell-1}d_\ell}\,\ip{\omega_i^{(\ell)}}{\omega_j^{(\ell)}},
\]
and by concatenating $\omega_i=(\omega_i^{(1)},\dots,\omega_i^{(L)})\in\RR^{m}$ we obtain
\[
  \sum_{\ell=1}^L K^{(\ell)} = \frac{1}{m}\Omega\Omega^\top,
\]
where $\Omega\in\RR^{n\times m}$ has rows $\omega_i$.

\section{Main result}
\label{sec:main}

\begin{thm}[Deep quadratic-scaling NTK bulk law]\label{thm:deep-main}
Assume \Cref{ass:scaling} and \Cref{ass:activation}.
Assume additionally that the deep replacement/Gaussian-equivalence step holds at each layer in the precise sense of the Benigni--Paquette resolvent replacement strategy, with deterministic $O(1/d_{\ell-1})$ fourth-order tensor deformations controlled by finite-width cumulant expansions.
Then there exists a deterministic probability measure $\mu_{\mathrm{deep}}$ on $[0,\infty)$ such that the empirical spectral distribution of $\Theta^{(L)}$ converges weakly in probability to
\[
  \mu_{\mathrm{MP}}^{\gamma}\boxtimes \mu_{\mathrm{deep}},
\]
where $\mu_{\mathrm{MP}}^{\gamma}$ is the Marchenko--Pastur law of shape $\gamma$.
\end{thm}

\section{Proof overview and dependency graph}
The proof follows the Benigni--Paquette dependency graph layerwise:
\begin{enumerate}
  \item Concentration: control quadratic forms in the deep features (Hanson--Wright type bounds).
  \item Tensor reduction: compute the conditional covariance tensor $T$ of $\omega_i$ and show it has a deterministic limiting ESD.
  \item Resolvent method: apply a Bai--Zhou type MP-map theorem to deduce the limiting ESD of $\frac{1}{m}\Omega\Omega^\top$.
\end{enumerate}
The only new step is the deep Gaussian-equivalence: we must replace the deep features entering $\omega_i$ by Gaussian surrogates with matching covariance and controlled fourth cumulants; finite-width Feynman-diagram cumulants provide the required $O(1/d)$ control.

\section{Appendix (outline)}
The appendices in a full version include:
(i) concentration lemmas for deep features,
(ii) a layerwise replacement proof mirroring BP,
(iii) a Bai--Zhou MP-map statement and verification of its hypotheses,
(iv) analysis of the leading covariance operator $Q$ and its spectrum in special cases.

\bibliography{references}

\end{document}

